\chapter{MongoDB Aggregation}

\section{The Pipeline Concept}

An aggregation pipeline runs documents through a sequence of stages, each
transforming the previous stage's output. Unlike \texttt{.find()}, it can
reshape, group, and compute across documents --- ideal for dashboard stats.

\begin{diagrambox}[Aggregation Pipeline]
\begin{verbatim}
Task collection
  $match   --> filter documents
  $group   --> bucket + compute per-bucket values
  $sort    --> order grouped results
  $project --> reshape output fields
  Final result array
\end{verbatim}
\end{diagrambox}

\section{Dashboard Stats Endpoint}

\texttt{GET /api/tasks/stats} needs total/completed/pending counts and a
per-user breakdown, via two \texttt{\$group} stages in one pipeline.

\begin{lstlisting}[language=JavaScript]
// backend/src/controllers/task.controller.js
const Task = require("../models/Task");

const getTaskStats = async (req, res, next) => {
  try {
    const statusCounts = await Task.aggregate([
      // Stage 1: only this user's tasks (skip if stats are app-wide)
      { $match: { owner: req.user.id } },

      // Stage 2: group by status, counting documents in each group
      {
        $group: {
          _id: "$status",
          count: { $sum: 1 },
        },
      },

      // Stage 3: alphabetical by status for a stable response order
      { $sort: { _id: 1 } },
    ]);

    // Reshape [{ _id: "completed", count: 5 }, ...] into a flat object
    const summary = { total: 0, completed: 0, pending: 0, "in-progress": 0 };
    statusCounts.forEach(({ _id, count }) => {
      summary[_id] = count;
      summary.total += count;
    });

    const perUser = await Task.aggregate([
      {
        $group: {
          _id: "$owner",
          taskCount: { $sum: 1 },
        },
      },
      // Join back to the users collection to attach a readable name
      {
        $lookup: {
          from: "users",
          localField: "_id",
          foreignField: "_id",
          as: "user",
        },
      },
      { $unwind: "$user" },
      {
        $project: {
          _id: 0,
          userId: "$_id",
          name: "$user.name",
          taskCount: 1,
        },
      },
      { $sort: { taskCount: -1 } },
    ]);

    res.status(200).json({ summary, perUser });
  } catch (error) {
    next(error);
  }
};

module.exports = { getTaskStats };
\end{lstlisting}

\begin{notebox}[NOTE]
\texttt{\$lookup} joins another collection --- here pulling each owner's
\texttt{name} from \texttt{users}. It's the aggregation-native equivalent of
\texttt{.populate()}, which isn't available inside a pipeline.
\end{notebox}

\subsection{Example Output}

\begin{lstlisting}[language=JavaScript]
{
  "summary": {
    "total": 45,
    "completed": 20,
    "pending": 15,
    "in-progress": 10
  },
  "perUser": [
    { "userId": "64f...a1", "name": "Asha Patel", "taskCount": 18 },
    { "userId": "64f...b2", "name": "Ravi Kumar", "taskCount": 15 },
    { "userId": "64f...c3", "name": "Meera Nair", "taskCount": 12 }
  ]
}
\end{lstlisting}

\section{Route Wiring}

\begin{lstlisting}[language=JavaScript]
// backend/src/routes/task.routes.js
router.get("/stats", protect, getTaskStats);
\end{lstlisting}

\begin{tipbox}[TIP]
Put \texttt{\$match} stages as early as possible --- they narrow the working
set before the costlier \texttt{\$group}/\texttt{\$lookup} stages run.
\end{tipbox}
